\setcounter{section}{6}






  \zwischenueberschrift{Übungsaufgaben}

\inputaufgabe
{}
{





Es sei

\mavergleichskette
{\vergleichskette
{ F(x)
}
{ \in }{ K[x]
}
{  }{
}
{    }{
}
{   }{
}
}
{}{}{}
ein Polynom in einer Variablen über einem
\definitionsverweis {Körper}{}{}
$K$. Parametrisiere den
\definitionsverweis {Graphen}{}{}
zu $F$ durch Polynome.





}
{} {}

\inputaufgabe
{}
{





Bestimme für die parametrisierte Kurve

\mathdisp {x= -3t^2+4t-2 \text{   und } y=2t^2+5t-3} {  }

eine Kurvengleichung.





}
{} {}

\inputaufgabegibtloesung
{}
{





Es sei $K$ ein
\definitionsverweis {Körper}{}{.}
Betrachte die durch
\maabbeledisp {} { {\mathbb A}^{1}_{K}  } { {\mathbb A}^{2}_{K}
} { t } {  \left(  t+t^2   , \,   t^3                   \right) = (x,y)
} {,}
definierte Parametrisierung. Bestimme eine
\zusatzklammer {nichttriviale} {} {}
algebraische Gleichung, die  für alle Bildpunkte dieser Abbildung erfüllt ist. Man gebe auch einen Punkt in der affinen Ebene an, der nicht auf der Bildkurve liegt.





}
{} {}

\inputaufgabegibtloesung
{}
{





Es sei $K$ ein Körper. Betrachte die durch
\maabbeledisp {} { {\mathbb A}^{1}_{ K } } { {\mathbb A}^{2}_{ K }
} { t } { \left(  t^2+1   , \,   t^3-t                   \right)
} {}
definierte polynomiale Abbildung. Bestimme eine
\zusatzklammer {nichttriviale} {} {}
algebraische Gleichung, die für alle Bildpunkte dieser Abbildung erfüllt ist.





}
{} {}

Das in der folgenden Aufgabe beschriebene Phänomen kann über einem algebraisch abgeschlossenen Körper nicht auftreten.

\inputaufgabe
{}
{





Man gebe ein Beispiel für eine polynomiale Parametrisierung
\maabbdisp {\varphi} { {\mathbb A}^{1}_{\R} } { C  \subset {\mathbb A}^{2}_{\R}
} {}
an, wobei $C$ der Zariski-Abschluss des Bildes sein soll und wobei unendlich viele Punkte von $C$ nicht im Bild liegen.





}
{} {}

\inputaufgabe
{}
{





Zeige, dass in der Situation von

Beispiel 6.3

die angegebene Parametrisierung der Kurve für

\mavergleichskette
{\vergleichskette
{ K
}
{ = }{ \R
}
{  }{
}
{    }{
}
{   }{
}
}
{}{}{}
surjektiv ist.





}
{} {}

\inputaufgabe
{}
{





Es seien

\mavergleichskette
{\vergleichskette
{ P_1,P_2,P_3
}
{ \in }{ K[S,T]
}
{  }{
}
{    }{
}
{   }{
}
}
{}{}{}
drei
\definitionsverweis {homogene Polynome}{}{}
vom gleichen Grad. Zeige, dass es ein homogenes Polynom

\mathbed {F \in K[X,Y,Z]} {}
 {F \neq 0} {}
 {} {} {} {,}
mit

\mavergleichskettedisp
{\vergleichskette
{ F(P_1,P_2,P_3)
}
{ =} { 0
}
{ } {
}
{ } {
}
{ } {
}
}
{}{}{}
gibt.





}
{} {}

\inputaufgabegibtloesung
{}
{





Gibt es ein
\definitionsverweis {homogenes Polynom}{}{}

\mavergleichskette
{\vergleichskette
{ F
}
{ \in }{ K[X,Y,Z]
}
{  }{
}
{    }{
}
{   }{
}
}
{}{}{}
\zusatzklammer {
\mavergleichskettek
{\vergleichskettek
{ F
}
{ \neq }{ 0
}
{  }{
}
{    }{
}
{   }{
}
}
{}{}{}} {} {}
mit

\mavergleichskette
{\vergleichskette
{ F(S,T,ST)
}
{ = }{ 0
}
{  }{
}
{    }{
}
{   }{
}
}
{}{}{?}





}
{} {}

\inputaufgabegibtloesung
{}
{





Bestimme für die Abbildung
\maabbeledisp {} { {\mathbb A}^{1}_{ K } \setminus \{0\}} { {\mathbb A}^{2}_{ K }
} { t } { \left(  { \frac{  t^2+1 }{ t  } }   , \,   { \frac{  t+1  }{ t } }                   \right)
} {,}
eine algebraische Gleichung der Bildkurve. Führe eine Probe durch.





}
{} {}

\inputaufgabe
{}
{





Bestimme für das Bild der Abbildung
\maabbeledisp {} { {\mathbb A}^{1}_{ K } \setminus \{-1,0,1\}} { {\mathbb A}^{2}_{ K }
} { t } { \left(  { \frac{   t  }{ t^2-1 } }   , \,   { \frac{   1  }{ t  } }                   \right)
} {}
eine nichttriviale algebraische Gleichung.





}
{} {}

In den folgenden Aufgaben wird auf den Begriff der differenzierbaren Kurve Bezug genommen, wie er in der Analysis 2 behandelt wird.

\inputaufgabe
{}
{





Man gebe ein Beispiel für eine nicht-algebraische
\zusatzklammer {nicht polynomiale} {} {}
\definitionsverweis {differenzierbare Kurve}{}{}
im $\R^2$, deren Bild aber mit einer
\definitionsverweis {algebraischen Kurve}{}{}
übereinstimmt.





}
{} {}

\inputaufgabe
{}
{





Zeige, dass die
\zusatzklammer {Bahn der} {} {}
archimedische Spirale
\maabbeledisp {f} {\R_{\geq 0}} {\R^2
} { t } { \left( t \cos  t   , \,  t \sin  t                   \right)
} {,}
nicht
\definitionsverweis {algebraisch}{}{}
ist.





}
{} {}

\inputaufgabe
{}
{





Es sei

\mavergleichskette
{\vergleichskette
{C
}
{ = }{ V(F)
}
{ \subset }{ {\mathbb A}^{2}_{\R}
}
{    }{
}
{   }{
}
}
{}{}{}
eine
\definitionsverweis {rational-parametrisierbare Kurve}{}{.}
Zeige, dass es auch nicht-algebraische
\definitionsverweis {differenzierbare}{}{}
Parametrisierungen der Kurve gibt.





}
{} {}

\inputaufgabe
{}
{





Man gebe eine
\definitionsverweis {differenzierbare Kurve}{}{}
\maabbdisp {f} {\R} {\R^2
} {}
an, deren
\definitionsverweis {Bild}{}{}
genau das
\definitionsverweis {Achsenkreuz}{}{}
ist.





}
{} {}

\inputaufgabe
{}
{





Es sei
\maabb {\varphi} { {\mathbb A}^{2}_{ K }  } { {\mathbb A}^{2}_{ K }
} {}
eine
\definitionsverweis {polynomiale Abbildung}{}{}
und sei $C$ eine ebene
\definitionsverweis {rationale}{}{}
Kurve. Es sei ferner vorausgesetzt, dass $C$ durch $\varphi$ nicht auf einen einzigen Punkt abgebildet wird. Zeige, dass dann $\overline{\varphi(C)}$ ebenfalls eine rationale Kurve ist.





}
{} {}

\inputaufgabe
{}
{





Es sei $K$ ein
\definitionsverweis {algebraisch abgeschlossener Körper}{}{}
der
\definitionsverweis {Charakteristik}{}{}
$0$. Es sei eine
\definitionsverweis {polynomiale Abbildung}{}{}
der Form
\maabbeledisp {\varphi} { {\mathbb A}^{1}_{K} } { {\mathbb A}^{2}_{K}
} { t } {   \left(  t^2   , \,   \psi(t)                   \right)
} {,}
gegeben
\zusatzklammer {mit

\mavergleichskettek
{\vergleichskettek
{ \psi(t)
}
{ \in }{ K[t]
}
{  }{
}
{    }{
}
{   }{
}
}
{}{}{}} {} {.}
Zeige, dass $\varphi$ genau dann injektiv ist, wenn $\psi$ die Form hat

\mavergleichskettedisp
{\vergleichskette
{ \psi(t)
}
{ =} { at^n+\theta(t)
}
{ } {
}
{ } {
}
{ } {
}
}
{}{}{}
mit

\mavergleichskette
{\vergleichskette
{ a
}
{ \neq }{ 0
}
{  }{
}
{    }{
}
{   }{
}
}
{}{}{,}
$n$ ungerade und $\theta(t)$ ein Polynom, in dem nur geradzahlige Exponenten auftreten.





}
{} {}

\inputaufgabegibtloesung
{}
{





Es sei

\mavergleichskette
{\vergleichskette
{ p
}
{ \in }{ [0,1]
}
{  }{
}
{    }{
}
{   }{
}
}
{}{}{}
die Wahrscheinlichkeit, mit der ein bestimmtes Ereignis $A$  bei der Durchführung eines Experiments eintritt, und entsprechend sei $1-p$ die Wahrscheinlichkeit, dass es nicht eintritt
\zusatzklammer {$\neg A$} {} {.}
Das Experiment werde zweimal
\definitionsverweis {unabhängig}{}{}
voneinander durchgeführt. Die möglichen Gesamtereignisse
\mathl{(A,A), (A, \neg A), (\neg A,A) , (\neg A, \neg A)}{} haben dann eine von $p$ abhängige Wahrscheinlichkeit. Diese Abhängigkeit fassen wir als die
\definitionsverweis {polynomiale Abbildung}{}{}
\maabbeledisp {\varphi} { {\mathbb A}^{1}_{ K }  } { { {\mathbb A}_{  K  }^{  4   } }
} { p } { \left(  p^2  , \,   p(1-p)  , \,   (1-p)p , \,   (1-p)^2                \right) = \left(  x  , \,   y  , \,   z , \,   w                \right)
} {,}
auf.
\aufzaehlungzweiabc{Zeige, dass die Abbildung
\definitionsverweis {injektiv}{}{}
ist.
}{Beschreibe das
\definitionsverweis {Bild}{}{}
der Abbildung vollständig durch polynomiale Gleichungen.
}





}
{} {}

\inputaufgabegibtloesung
{}
{





Es sei

\mavergleichskette
{\vergleichskette
{ p
}
{ \in }{ [0,1]
}
{  }{
}
{    }{
}
{   }{
}
}
{}{}{}
die Wahrscheinlichkeit, mit der ein bestimmtes Ereignis $A$ bei der Durchführung eines ersten Experiments eintritt und  sei

\mavergleichskette
{\vergleichskette
{ q
}
{ \in }{ [0,1]
}
{  }{
}
{    }{
}
{   }{
}
}
{}{}{}
die Wahrscheinlichkeit, mit der ein bestimmtes Ereignis $B$ bei der Durchführung eines zweiten Experiments eintritt. Die beiden Experimente werden
\definitionsverweis {unabhängig}{}{}
voneinander durchgeführt. Die möglichen Gesamtereignisse
\mathl{(A,B), (A, \neg B), (\neg A,B) , (\neg A, \neg B)}{} haben dann eine von $p$ und $q$ abhängige Wahrscheinlichkeit. Diese Abhängigkeit fassen wir als die
\definitionsverweis {polynomiale Abbildung}{}{}
\maabbelementzeiledisplay {\varphi} { {\mathbb A}^{2}_{ K }  } { { {\mathbb A}_{  K  }^{  4   } }
} {(p,q)} { \left(  pq  , \,   p(1-q)  , \,   (1-p)q , \,   (1-p)(1-q)                \right) = \left(  x  , \,   y  , \,   z , \,   w                \right)
} {,}
auf.
\aufzaehlungzweiabc{Zeige, dass die Abbildung
\definitionsverweis {injektiv}{}{}
ist.
}{Beschreibe das
\definitionsverweis {Bild}{}{}
der Abbildung vollständig durch polynomiale Gleichungen.
}





}
{} {}

\inputaufgabe
{}
{





Wir betrachten die Abbildung
\maabbeledisp {} { {\mathbb A}^{2}_{ K } \supseteq D(s) } { { {\mathbb A}_{ K }^{  3   } }
} { (s,t) } { \left(  s  , \,   { \frac{  t^2 }{ s } }  , \,   t                 \right) = (x,y,z)
} {.}
Bestimme eine algebraische Gleichung $F$ für das Bild. Untersuche die Abbildung auf Injektivität und Surjektivität (als Abbildung nach $V(F)$). Vergleiche diese Abbildung mit den in

Aufgabe 6.29

diskutierten Abbildungen.





}
{} {}

\inputaufgabe
{}
{





Es sei $K$ ein
\definitionsverweis {Körper}{}{}
und
\mathl{K[X_1 , \ldots ,  X_n]}{} der
\definitionsverweis {Polynomring}{}{}
über $K$ in $n$ Variablen und
\mathl{K[X_1 , \ldots , X_n,Z]}{} der Polynomring in
\mathl{n+1}{} Variablen. Zu

\mavergleichskette
{\vergleichskette
{ F
}
{ \in }{ K[X_1 , \ldots , X_n]
}
{  }{
}
{    }{
}
{   }{
}
}
{}{}{}
sei

\mavergleichskette
{\vergleichskette
{  \hat{ F  }
}
{ \in }{ K[X_1 , \ldots , X_n,Z]
}
{  }{
}
{    }{
}
{   }{
}
}
{}{}{}
die
\definitionsverweis {Homogenisierung}{}{}
\zusatzklammer {bezüglich $Z$} {} {}
und zu

\mavergleichskette
{\vergleichskette
{ G
}
{ \in }{ K[X_1 , \ldots , X_n,Z]
}
{  }{
}
{    }{
}
{   }{
}
}
{}{}{}
sei
\mathl{\tilde{ G }}{} die
\zusatzklammer {durch
\mathl{Z \mapsto 1}{} gegebene} {} {}
\definitionsverweis {Dehomogenisierung}{}{}
von $G$. Zeige, dass

\mavergleichskette
{\vergleichskette
{  \tilde{  \hat{  F   }  }
}
{ = }{ F
}
{  }{
}
{    }{
}
{   }{
}
}
{}{}{,}
aber nicht

\mavergleichskette
{\vergleichskette
{ \hat{  \tilde{  G  }   }
}
{ = }{ G
}
{  }{
}
{    }{
}
{   }{
}
}
{}{}{}
gelten muss.





}
{} {}

\inputaufgabegibtloesung
{}
{





Es sei $K$ ein
\definitionsverweis {Körper}{}{} und
\mathl{K[X_1 , \ldots ,  X_n,Z]}{} der
\definitionsverweis {Polynomring}{}{} über $K$ in
\mathl{n+1}{} Variablen. Es seien

\mavergleichskette
{\vergleichskette
{ G,H
}
{ \in }{ K[X_1 , \ldots , X_n,Z]
}
{  }{
}
{    }{
}
{   }{
}
}
{}{}{}
\definitionsverweis {homogene Polynome}{}{}
vom gleichen Grad. Für die
\definitionsverweis {Dehomogenisierungen}{}{}
\zusatzklammer {bezüglich $Z$} {} {}
gelte

\mavergleichskette
{\vergleichskette
{ \tilde{  G  }
}
{ = }{ \tilde{  H  }
}
{  }{
}
{    }{
}
{   }{
}
}
{}{}{.}
Zeige, dass dann

\mavergleichskette
{\vergleichskette
{ G
}
{ = }{ H
}
{  }{
}
{    }{
}
{   }{
}
}
{}{}{}
ist.





}
{} {}

\inputaufgabegibtloesung
{}
{





Es sei $K$ ein
\definitionsverweis {Körper}{}{}
und
\mathl{K[X_1 , \ldots ,  X_n]}{} der
\definitionsverweis {Polynomring}{}{}
über $K$ in $n$ Variablen und
\mathl{K[X_1 , \ldots , X_n,Z]}{} der Polynomring in
\mathl{n+1}{} Variablen. Zeige, dass die
\definitionsverweis {Homogenisierung}{}{}
\zusatzklammer {bezüglich $Z$} {} {}
mit der Multiplikation verträglich ist.





}
{} {}

\inputaufgabe
{}
{





Es sei $K$ ein
\definitionsverweis {Körper}{}{}
und
\mathl{K[X_1 , \ldots ,  X_n]}{} der
\definitionsverweis {Polynomring}{}{}
über $K$ in $n$ Variablen und
\mathl{K[X_1 , \ldots , X_n,Z]}{} der Polynomring in
\mathl{n+1}{} Variablen. Beschreibe die
\definitionsverweis {Dehomogenisierung}{}{}
\zusatzklammer {bezüglich $Z$} {} {}
als einen
\definitionsverweis {Einsetzungshomomorphismus}{}{.}





}
{} {}

\inputaufgabe
{}
{





Formuliere und beweise eine Division mit Rest-Aussage für homogene Polynome in zwei Variablen über einem Körper.





}
{} {}

\inputaufgabegibtloesung
{}
{





Es seien

\mavergleichskettedisp
{\vergleichskette
{F
}
{ =} { X^4+9X^3Y+7X^2Y^2+XY^3+8Y^4
}
{ } {
}
{ } {
}
{ } {
}
}
{}{}{}
und

\mavergleichskettedisp
{\vergleichskette
{G
}
{ =} { X^3 +5X^2Y
}
{ } {
}
{ } {
}
{ } {
}
}
{}{}{.}
Finde
\definitionsverweis {homogene Polynome}{}{}

\mathl{Q,R  \in \Q[X,Y]}{,}
\mathl{Q \neq 0}{,} mit

\mavergleichskettedisp
{\vergleichskette
{F
}
{ =} { GQ+R
}
{ } {
}
{ } {
}
{ } {
}
}
{}{}{.}





}
{} {}





  \zwischenueberschrift{Aufgaben zum Abgeben}

\inputaufgabe
{3}
{





Bestimme den Flächeninhalt der durch die reelle Kurve

\mavergleichskettedisp
{\vergleichskette
{V(y^2-x^3-x^2)
}
{ \subset} { {\mathbb A}^{2}_{\R}
}
{ } {
}
{ } {
}
{ } {
}
}
{}{}{}
eingeschlossenen Schlaufe.





}
{} {}

Die folgenden Aufgaben erfordert eventuell den Einsatz eines Computers.

\inputaufgabe
{6}
{





Bestimme für die Abbildung
\maabbeledisp {} { {\mathbb A}^{1}_{ K } } { {\mathbb A}^{2}_{ K }
} { t } {   \left(  t^2+t^3   , \,   2t^2-t^4                   \right)
} {,}
eine algebraische Gleichung der Bildkurve.





}
{} {}

\inputaufgabe
{5}
{





Bestimme für die Abbildung
\maabbeledisp {} { {\mathbb A}^{1}_{ K } \setminus \{0\}} { {\mathbb A}^{2}_{ K }
} { t } { \left(  { \frac{  t^2-1 }{ t  } }   , \,   { \frac{  t+3  }{ t } }                   \right)
} {,}
eine algebraische Gleichung der Bildkurve. Führe eine Probe durch.





}
{} {}

\inputaufgabe
{5}
{





Wir betrachten die beiden Abbildungen

\mathdisp {(s,t) \longmapsto  \left(  s^2   , \,   t^2  , \,   st                 \right) = (x,y,z) \text{ und } (s,t) \longmapsto \left(  s   , \,   st^2  , \,   st                 \right) = (x,y,z)} { . }

Zeige, dass das Bild der beiden Abbildungen die gleiche algebraische Gleichung $F$ erfüllt. Untersuche die Abbildungen auf Injektivität und Surjektivität
\zusatzklammer {als Abbildung nach $V(F)$} {} {.}
Welche Abbildung liefert eine \anfuehrung{bessere}{} Beschreibung von $V(F)$?





}
{} {}

\inputaufgabe
{4}
{





Es sei $K$ ein
\definitionsverweis {Körper}{}{}
und

\mavergleichskette
{\vergleichskette
{ F
}
{ \in }{ K[X,Y]
}
{  }{
}
{    }{
}
{   }{
}
}
{}{}{}
ein
\definitionsverweis {irreduzibles}{}{}
Polynom. Die
\definitionsverweis {Nullstellenmenge}{}{}

\mathl{V(F)}{} sei unendlich. Zeige, dass dann
\mathl{V(F)}{} eine
\definitionsverweis {irreduzible}{}{}
\definitionsverweis {affin-algebraische}{}{}
Menge ist.





}
{Man gebe auch ein Beispiel, dass diese Aussage in drei Variablen falsch ist.} {}
